\chapter{Discussion}
\label{ch:discussion}

\section{Designing for sensor loss was worth it, and not for the reason expected}

The premise was that a body-worn control arm routing analogue signals through
turning joints will lose channels, so losing them has to be survivable. That
premise held. Seven of fourteen channels currently fail the health test, and
the failures are progressive rather than sudden.

What the architecture actually bought is narrower than it first appears. It
did not keep the system working: with the channels in their present state,
direct teleoperation is not usable, and no amount of graceful degradation
changes that. What it bought was the ability to say so precisely. The system
names which channels it no longer believes, states what capability that costs,
and refuses rather than inventing a position when too little is left. The
alternative, which is what the system did before, was to keep commanding the
arm from readings that were confidently wrong.

There is an uncomfortable corollary. A system that degrades gracefully also
degrades quietly. Six of the fourteen channels failed over a period during
which nothing stopped working in a way anybody noticed, and the failure was
found by measuring rather than by using. Any design like this needs the health
report to be visible in the ordinary course of a session, not available on
request, which is why the channel state now appears in the operating window
rather than in a log.

\section{The reduced-control premise, assessed}

The second premise was that commanding fourteen joints is a usability problem,
so the system should take position from a few well-behaved sensors and leave
the rest to the machine or to nobody.

Taking position from a small number of channels worked. Splitting it into
height, direction and distance, and taking each from whichever sensor sees it
most directly, is the reason a system with half its sensors failing still
produces a usable command on one arm.

Pinning the wrist did not work as well as expected, and the cost was not
measured until late. Fixing the wrist orientation on a seven-jointed arm uses
up the spare joint, and the spare joint is precisely the one that swings the
elbow away from the wearer while the hand stays put. Measured at the point
where the work actually happens, pinning the wrist reduced the average
distance the arm could travel from \SI{0.30}{\metre} to \SI{0.065}{\metre}, a
factor of seven. Allowing the wrist to vary within a \SI{15}{\degree} cone
recovers most of it and costs nothing measurable in accuracy. This should have
been measured a month earlier than it was.

\section{The workspace result, and what it costs the study}

The two arms cannot both reach any point in front of the wearer. That result
was checked four different ways precisely because it invalidates part of the
study, and it survived all four.

The consequences are concrete. Any task requiring the two hands to meet is
impossible, which removes the passing of an object from one arm to the other.
Anything requiring both arms on the same object requires the object to be wide
enough to span the dead band, which is where the two-arm carry task ran into
the gripper's opening. And the innermost column of the working area sits
\SI{325}{\milli\metre} from the body's centre line, so a task described as
happening in front of the wearer actually happens off to one side.

It is worth being clear about what kind of result this is. It is not a
limitation of the control system, and no improvement to the software affects
it. It is a consequence of where two arms of this length are bolted to a
person of this size. The measured mount change would improve it substantially
and has not been applied, because applying it means re-deriving the resting
pose, the approach direction and every clearance figure in this thesis. That
is a session's work, not an afternoon's, and it should be done as one
deliberate pass with everything re-measured afterwards.

\section{The safety layer is honest about the wrong thing}

The safety layer is thorough about distance and silent about time.

Fourteen faults are injected and handled, the margin is measured
geometrically against the parts of the arm that actually pass the wearer, and
every mechanism that can stop the arm announces itself by name. All of that is
sound and most of it was built after a specific failure showed why it was
needed.

What none of it addresses is that the margin is a distance and the wearer can
move. In five of six modelled limb movements, the limb travels further than
the entire \SI{150}{\milli\metre} margin before the guard hears about it. The
system is protecting a stationary person from a moving robot, and the
arrangement it is deployed in has a moving person as well.

Three things would help and none is done. Inflating the margin by the delay
times an assumed speed would at least make the number mean something, and it
is a one-parameter change. Cutting the delay is mostly achievable, because
most of the \SI{562}{\milli\second} is a rate limit rather than the tracking
itself. And slowing the arm as it approaches the wearer would couple two
quantities the system already has and currently does not connect. The first is
the one to do first, because it changes what the margin claims rather than
merely improving it.

\section{Measurement was most of the work}

The single most striking number in this project is not about the robot. It is
that a surprising result turned out to be a fault in the measurement
seventeen times.

That has three consequences for how this thesis should be read. Any figure
here that has not been produced by a program with a known-answer test behind it
should be treated as provisional, and the text says which those are. Any
figure that was produced once should be treated as provisional too, because
the solver used here restarts from a random guess and a single call is a coin
flip. And the absence of a measurement is reported rather than filled in,
which is why several sentences in \cref{ch:results} end in a statement that
something has never been measured.

The habit is not free. Clearing an instrument before believing it roughly
doubles the time a measurement takes. Set against that, several of the
instrument faults found here would have propagated into published numbers, and
two of them into the safety case.

\section{Limitations}

\emph{Nothing has run on a moving arm under load.} The simulator used here
returns zero for every force and torque, so anything derived from a force
reading is undefined on this platform rather than merely unmeasured. That is
why there is no force feedback result in this thesis and no impedance
criterion in the evaluation.

\emph{The delay figures are lower bounds.} They are measured within one
machine, over shared memory. The physical arm adds a network round trip per
command, and that has not been characterised.

\emph{The mapping results depend on a broken master.} The sideways mapping
problem is reported as geometric, and the argument for that is a
rotation-invariant measurement rather than an inference from the failing
channels. But the accuracy figures for the working channels come from an arm
whose other half does not work, and they should be re-measured after repair.

\emph{One person built and measured all of it.} Every design decision was
taken by the person who then measured its consequences. The instrument-checking
habit was adopted partly as a defence against that, and it is not a complete
one.

\section{What stands between this and the study}

In the order it should be done.

\begin{enumerate}
  \item \textbf{Repair the master arm's wiring.} Two of the fourteen channels
    are dead and six are incoherent, and the failures cluster on one arm. Until
    that is fixed there is no direct teleoperation baseline, and without a
    baseline there is no comparison to make. This is a soldering job.
  \item \textbf{Measure where the wrist camera actually is.} It is the largest
    unknown in the positioning budget and the reason no grasp has landed from
    an open-loop plan. It needs the hardware and a calibration target.
  \item \textbf{Make the clearance check measure the whole arm.} At present it
    samples the centres of joints, so it can return a comfortable number for a
    pose in which the arm is touching the mannequin. This is the outstanding
    safety defect and it should be fixed before anybody wears the rig.
  \item \textbf{Put the delay into the margin.} One parameter, and it turns a
    distance that means little into one that means something.
  \item \textbf{Apply the mount change and re-measure everything.} A day's work
    and a large improvement in what the arms can reach.
  \item \textbf{Then submit for ethical approval}, with a baseline that works
    and a safety case whose numbers have been measured on the arm rather than
    in simulation.
\end{enumerate}
